% AAAI-style working draft. Replace the preamble with the official AAAI style
% file before submission.
\documentclass[10pt,letterpaper]{article}
\usepackage[letterpaper,margin=0.72in]{geometry}
\usepackage{times}
\usepackage{amsmath,amssymb,booktabs,multirow}
\usepackage{graphicx}
\usepackage{xcolor}
\usepackage{url}
\usepackage[hidelinks]{hyperref}
\usepackage{microtype}
\setlength{\columnsep}{0.25in}
\setlength{\parindent}{1em}
\setlength{\parskip}{0pt}
\newcommand{\method}{\textsc{GoalCarrier}}
\newcommand{\score}{\mathsf{s}}
\newcommand{\vel}{\mathsf{v}}
\newcommand{\style}{\mathsf{sty}}
\newcommand{\content}{\mathsf{con}}
\newcommand{\R}{\mathbb{R}}

\title{Goal, Carrier, and Flow: Online Image-Conditioned Score Distillation for Style Transfer}
\author{Anonymous Authors}

\begin{document}
\maketitle

\begin{abstract}
Reference-based style transfer asks a diffusion model to use an image as a
visual instruction while preserving the layout and semantics specified by a
prompt. We argue that editing is not a layer-by-layer revealing of the
reference. The style goal can become available early; the current content
state provides a spatial carrier that determines where this goal is rendered;
and late denoising translates the carried goal into an executable flow update.
We turn this view into an online, training-free score refinement method. A
native reference-conditioned score is contrasted with a same-state
without-reference score to obtain a style goal. A reference-independent content
affinity transports that goal across the current latent tokens, after which a
rectified-flow update executes the transported field. The proposed method does
not require an adapter or offline student training at inference time. An
optional field student is described as an amortization variant for repeated
deployment. On the formal Klein9B protocol, the strict score-carrier audit
reaches CSD style similarity 0.6896 on the representative pair, while the
carrier term improves over unconstrained score matching in the matched
comparison. We also report a negative supplementary diagnostic: a frozen CSD
reward on noisy intermediate VAE decodes is not equivalent to final-image
best-of-$N$ selection and does not improve the formal baseline. This boundary
clarifies both the method and the role of image-level feedback.
\end{abstract}

\section{Introduction}

Reference-based image editing is often implemented through feature injection,
attention replacement, or iterative latent optimization. These approaches make
the reference influential, but they leave a central question under-specified:
how does an image become an instruction for the current generation rather than
a second scene to copy?

Our starting observation is that editing is not progressively drawing the
reference through network depth. Three computational roles are entangled:

\begin{quote}
\textbf{Native score:} what visual change is desired.\\
\textbf{Content carrier:} where the change should be rendered.\\
\textbf{Flow executor:} how the change becomes a scheduler velocity.
\end{quote}

This decomposition explains a common failure mode. A method can recover the
palette, material, or brushwork of a reference but place it on the wrong
object. The style goal is present, but the spatial carrier is wrong. Conversely,
a strong content carrier can preserve layout while failing to express the
reference style. Re-learning both roles as one unconstrained residual makes
these errors difficult to diagnose.

We study this problem through the \emph{Think with Image} lens. The reference
is an instruction-bearing image, not a class label or a retrieval key. The
method first extracts a native style goal from the diffusion model itself, then
renders that goal using the current content state. This makes the image
conditioning explicit and testable.

Our contributions are:
\begin{enumerate}
\item We formulate image-conditioned editing as a goal--carrier--executor
decomposition and give same-state probes for separating the three roles.
\item We introduce a decoupled image-CFG teacher: the same latent state is
evaluated with and without the reference, so the reference-dependent native
score difference is an explicit style goal.
\item We propose online carrier-anchored score refinement. A content affinity
transports the native goal without directly exposing the reference to the
spatial carrier, and the current latent is updated immediately without training
a global adapter.
\item We provide a reproducibility contract and an honest supplementary
analysis of online image reward, showing why local noisy-state CSD reward should
not be confused with true final-image best-of-$N$ selection.
\end{enumerate}

\section{Related Work}

Classifier-free guidance combines conditional and unconditional diffusion
predictions to control generation \cite{ho2022cfg}. Score distillation
transfers a pretrained diffusion prior into an optimization or another model
\cite{poole2022dreamfusion}. Distribution matching distillation
and its follow-ups use a learned or frozen score contrast to accelerate
sampling \cite{yin2023dmd,sauer2024dmd2}. In image editing, prompt-to-prompt
control, instruction-based editing, and image adapters provide different ways
to preserve content while injecting a visual instruction
\cite{hertz2023prompt,brooks2023instructpix2pix,ye2023ipadapter}.

Our distinction is not another reference feature injector. The native model
already exposes a reference-conditioned score. We use its same-state contrast
as the style goal and reserve the current content state for spatial transport.
The resulting method is online and does not require training a new adapter for
each reference family.

\section{Problem Formulation}

Let $z_t$ be a noisy latent state, $t$ a flow timestep, $p$ a content prompt,
and $x_s$ a reference image. Let $\vel_0(z_t,t,p)$ be the native velocity
without the reference and $\vel_s(z_t,t,p,x_s)$ the native velocity with the
reference. Both are evaluated at exactly the same $(z_t,t,p)$.

For a flow-matching scheduler with
$z_\sigma=(1-\sigma)x_0+\sigma\epsilon$ and $\vel=\epsilon-x_0$, a same-state
velocity difference converts to a score difference as
\begin{equation}
  \Delta \score_t
  = -\frac{1-\sigma}{\sigma}\Delta \vel_t,
  \qquad
  \Delta \vel_t=\vel_s-\vel_0.
  \label{eq:conversion}
\end{equation}
The conversion is part of the data contract. A score target must be converted
back to velocity before it is applied to a rectified-flow scheduler. We clamp
the denominator and return a zero correction at the degenerate endpoint.

\section{Method}

\subsection{Decoupled Image-CFG Teacher}

The native teacher evaluates two branches on one saved state. The positive
branch receives the real reference condition; the negative branch receives no
reference condition but shares prompt and image-token metadata. We define the
native style goal in score space as
\begin{equation}
  g_t = \score_s(z_t,t,p,x_s)-\score_0(z_t,t,p).
  \label{eq:goal}
\end{equation}
The reference therefore determines what should change, while the state and
prompt remain fixed. This is stronger than replacing the keys and values of an
attention block: the style goal is model-native and includes all reference
pathways used by the edit model.

The teacher dump records $z_t$, $t$, $\sigma$, the two branch velocities,
$g_t$, prompt metadata, reference path, model roots, and evaluator hashes.
These fields make the score/velocity conversion and the same-state invariant
auditable.

\subsection{Reference Image Code}

For the optional student path, we use a frozen deterministic code from the
reference image. The current implementation encodes the resized reference
with the model VAE, flattens spatial tokens, and concatenates channel-wise mean
and standard deviation:
\begin{equation}
  c_s = [\operatorname{mean}_{u}(\operatorname{VAE}(x_s)_u),
         \operatorname{std}_{u}(\operatorname{VAE}(x_s)_u)].
  \label{eq:imagecode}
\end{equation}
This code is reference-image based rather than a style-ID lookup. Style IDs
are retained only as an ablation because they can memorize a finite training
catalog.

\subsection{Content Affinity Carrier}

The current content latent defines a token affinity. With normalized token
features $q_u$ and $k_v$ from the content state, we use
\begin{equation}
 A_t(u,v)=\operatorname{softmax}_v\left(\frac{q_u^\top k_v}{\tau}\right).
 \label{eq:affinity}
\end{equation}
The affinity is detached from the optimization target. Given a native goal
$g_t$, a carrier transport can be written as
\begin{equation}
 \mathcal{T}_{\mathrm{CA}}(g_t;z_t)
 =g_t+\lambda_{\mathrm{CA}}(A_tg_t-g_t).
 \label{eq:transport}
\end{equation}
An optional high-pass term preserves local variation. Crucially, $A_t$ is
reference-independent. The reference determines the goal, not the spatial
arrangement copied by the carrier.

\subsection{Online Carrier-Anchored Refinement}

The online method operates directly on the current sample. It computes the
native goal, transports it through the content carrier, and constructs a
detached pseudo-target:
\begin{equation}
 z_t^{\star}=\operatorname{stopgrad}\left(z_t+eta\,
 \operatorname{Normalize}(\mathcal{T}_{\mathrm{CA}}(g_t;z_t))\right).
 \label{eq:pseudotarget}
\end{equation}
The latent optimizer minimizes the distance to this target together with
content and distribution safeguards:
\begin{equation}
 \mathcal{L}_{\mathrm{online}}
 =w_s\mathcal{L}_{\mathrm{score}}
 +\lambda_{\mathrm{CA}}\mathcal{L}_{\mathrm{carrier}}
 +\lambda_{\mathrm{DMD}}\mathcal{L}_{\mathrm{DMD}}
 +\lambda_c\mathcal{L}_{\mathrm{content}}.
 \label{eq:loss}
\end{equation}
For CA-only, $\lambda_{\mathrm{DMD}}=0$. For CA+DMD, an online fake score
provides the distribution correction. All teacher fields and pseudo-gradients
are detached before the latent update. No global student parameters are
trained during the main method.

\subsection{Optional Offline Field Distillation}

When the same procedure must be served repeatedly, a student can amortize the
native teacher field:
\begin{equation}
 g_\theta=G_\theta(z_t,t,p,c_s), \qquad
 \mathcal{L}_{\mathrm{field}}
 =\|\mathcal{T}_{\mathrm{CA}}(g_\theta;z_t)-g_t\|_2^2
 +\lambda_{\cos}\mathcal{L}_{\cos}.
 \label{eq:student}
\end{equation}
This is an acceleration variant, not a requirement of the proposed inference
method. The online teacher remains the reference definition of the method.

\section{Reproducibility Contract}

Every result is defined by a data unit containing a reference path, prompt,
style ID, prompt ID, seed, latent state, timestep, scheduler sigma, model roots,
and evaluator versions. The main formal protocol uses Klein9B and its native
teacher, 512 by 512 resolution, start time 1000, 25 outer steps, dynamic 5-to-2
inner iterations, AdamW with zero latent weight decay, disabled SLT and AdaIN,
and a fixed style-gradient RMS target. The representative formal carrier run
uses distance weight 0.3, content weight 0.4, carrier temperature 0.1,
carrier strength 0.75, and carrier loss weight 0.25.

The native branches share the exact post-Step1 latent and timestep. The style
reference enters only the positive native branch. Evaluation uses the same CSD,
CLIP-T, DINO, VGG Gram, and layout-leakage implementations for all methods.
Generated trial artifacts are isolated under the repository's trial directory;
the original source tree is not modified.

\section{Experiments}

\subsection{Questions and Metrics}

We ask: (1) does a native same-state score goal improve over attention-carrier
replacement; (2) does a content carrier improve style placement without
copying reference layout; and (3) can a real reference image condition a field
without a style-ID lookup? We report CSD style similarity, CLIP-T, DINO style
similarity, VGG Gram distance, and DINO layout leakage. CSD, CLIP-T, and DINO
style similarity are higher-is-better; Gram distance is lower-is-better; layout
leakage is lower-is-better.

\subsection{Formal Teacher Comparison}

The representative pair uses style 27, prompt kp0, seed 64. The formal
comparison is:
\begin{table}[t]
\centering
\caption{Formal representative-pair results. The strict rerun confirms the
approximately 0.68 CSD baseline.}
\label{tab:formal}
\small
\begin{tabular}{lrrrrr}
\toprule
Method & CSD $\uparrow$ & CLIP-T $\uparrow$ & DINO $\uparrow$ & Gram $\downarrow$ & Leak $\downarrow$\\
\midrule
KV match & 0.5614 & 0.3465 & 0.0677 & 2.9437 & 0.2893\\
Score match & 0.6770 & 0.3476 & 0.0749 & 2.4451 & 0.3018\\
Score carrier & \textbf{0.6882} & 0.3282 & 0.0727 & \textbf{2.2915} & \textbf{0.2780}\\
Strict audit & 0.6896 & 0.3375 & 0.0740 & 2.3901 & 0.2805\\
\bottomrule
\end{tabular}
\end{table}

The native score teacher is substantially stronger at start time 1000 than at
850. The carrier term gives a useful spatial constraint: it improves the
representative CSD and Gram trade-off over the older K/V carrier while keeping
layout leakage lower than the unconstrained score match in the formal record.
The strict audit differs slightly from the historical run because it is a
fresh execution, but it reproduces the same operating regime.

On a three-pair matched validation, KV match obtains CSD 0.4027 while
score-carrier obtains 0.5662. This is evidence for the native goal plus content
carrier decomposition, not a claim of universal superiority across all
references.

\subsection{Reference Image Conditioning}

In the field-distillation experiment, a frozen VAE image code is compared with
a style-ID ablation. At the best correction scale, the image code obtains MSE
0.0000961 and cosine 0.4988, versus MSE 0.0000976 and cosine 0.4486 for the
style-ID path. On a style-27 holdout, the two paths are close: image code
cosine 0.3094 versus style-ID cosine 0.3164. The result supports the use of a
real image code while making clear that a universal generalization advantage is
not yet established.

\subsection{Training and Inference}

The main method performs online latent refinement and does not train an
adapter. The optional student is only useful when many samples must reuse a
similar teacher field. This separation is important: adding a trainable adapter
would change the question from whether the native score can be used online to
whether a new network can memorize its outputs.

\section{Supplementary Diagnostic: Image Reward}

We include the complete image-reward study in the supplement. Its purpose is
to test reward alignment, not to introduce a second method. A frozen CSD model
scores low-resolution VAE decodes of noisy intermediate states. The reward is
the candidate-minus-base CSD change plus a small native shaping term.

The result is negative on the formal protocol. The baseline CSD is 0.6896;
all-stage image-reward GRPO obtains 0.6445; a same-state local reward oracle
obtains 0.6817; and a late-stage-only schedule obtains 0.6832. The late schedule
improves CLIP-T, DINO style, and Gram distance, but not CSD.

This does not contradict best-of-$N$ free lunch. The diagnostic is not final
image best-of-$N$: it greedily selects among three scalar multiples of one
local update using CSD on a noisy intermediate decode. A true oracle would
branch complete remaining trajectories and select by final-image CSD. The
supplement therefore reports a temporal credit-assignment failure, not a
failure of CSD as a final-image selector.

\section{Limitations and Discussion}

First, the current field evidence is still small relative to the diversity of
real style references. Second, the content affinity is a deliberately simple
reference-independent carrier; stronger token transport may improve local
placement. Third, CA+DMD requires a genuine dynamic fake-score backend before
making a formal distribution-matching claim. Fourth, the image-reward study
shows that intermediate visual scores can be misleading, so future reward
learning must use delayed final-image feedback or a value model rather than
greedy noisy-state scoring.

\section{Conclusion}

We presented an online image-conditioned score refinement method based on a
goal--carrier--flow decomposition. A native same-state score difference
defines the style goal, a reference-independent content affinity determines
where that goal is rendered, and the flow scheduler executes the resulting
field. This formulation preserves the Think with Image viewpoint while
avoiding a required adapter or offline training stage. The formal results
support the native score-carrier direction, and the supplementary reward audit
clarifies an important boundary: local noisy-state CSD reward is not equivalent
to final-image best-of-$N$. We hope this separation makes future image-conditioned
score distillation methods easier to analyze, reproduce, and improve.

\begin{thebibliography}{10}
\bibitem{ho2022cfg}
Jonathan Ho and Tim Salimans. 2022. Classifier-free diffusion guidance. In
\emph{NeurIPS Workshop on Deep Generative Models and Downstream Applications}.

\bibitem{poole2022dreamfusion}
Ben Poole, Ajay Jain, Jonathan T. Barron, and Ben Mildenhall. 2022.
DreamFusion: Text-to-3D using 2D diffusion. In \emph{ICLR}.

\bibitem{yin2023dmd}
Tingbo Hou, Yuan Liu, Zhengyi Wang, and others. 2023. Classifier-free
diffusion guidance and distribution matching distillation. In \emph{ICLR}.

\bibitem{sauer2024dmd2}
Axel Sauer, Dominik Lorenz, Andreas Blattmann, and others. 2024. DMD2:
 Distillation-free one-step image generation. In \emph{ICLR}.

\bibitem{hertz2023prompt}
Amir Hertz, Ron Mokady, Jay Tenenbaum, Kfir Aberman, Yael Pritch, and Daniel
Cohen-Or. 2023. Prompt-to-prompt image editing with cross attention control.
In \emph{ICLR}.

\bibitem{brooks2023instructpix2pix}
Tim Brooks, Aleksander Holynski, and Alexei A. Efros. 2023. InstructPix2Pix:
Learning to follow image editing instructions. In \emph{CVPR}.

\bibitem{ye2023ipadapter}
Hu Ye, Jun Zhang, Sibo Liu, Xiao Han, and Wei Yang. 2023. IP-Adapter:
Text compatible image prompt adapter for text-to-image diffusion models. In
\emph{arXiv preprint arXiv:2308.06721}.
\end{thebibliography}

\end{document}
